% ECONOMICS_PROBLEM: {"id": "H77-422", "title": "Fiscal federalism and intergovernmental incentives", "jel_code": "H77", "source_id": "OP-0048"}
% !TeX program = xelatex
\documentclass[11pt,a4paper]{article}
\usepackage[margin=23mm]{geometry}
\usepackage{fontspec}
\setmainfont{Latin Modern Roman}
\usepackage{amsmath,amssymb,mathtools}
\usepackage{xcolor,enumitem,booktabs,longtable,array,xurl,hyperref}
\definecolor{muted}{HTML}{53636C}
\hypersetup{colorlinks=true,urlcolor=blue,linkcolor=blue}
\setlength{\parindent}{0pt}
\setlength{\parskip}{5pt}
\newcommand{\R}{\mathbb R}
\newcommand{\E}{\mathbb E}
\newcommand{\N}{\mathbb N}
\newcommand{\ind}{\mathbf 1}
\newcommand{\Var}{\operatorname{Var}}
\newcommand{\Cov}{\operatorname{Cov}}
\newcommand{\supp}{\operatorname{supp}}
\DeclareMathOperator*{\argmin}{arg\,min}
\DeclareMathOperator*{\argmax}{arg\,max}
\newcommand{\entrymeta}[3]{{\sffamily\small\color{muted}\textbf{\texttt{#1}}\quad|\quad #2\par #3\par}}
\newcommand{\Problemref}[1]{\texttt{#1}}
\newcommand{\JELref}[2]{\texttt{#2} (JEL \texttt{#1})}
\begin{document}
\section*{Fiscal federalism and intergovernmental incentives}
\textbf{Problem ID:} \texttt{H77-422}\quad \textbf{JEL:} \texttt{H77}\par
% BEGIN ECONOMICS STATEMENT
\label{problem:OP-0048}
\label{atlas:prob:P8}
\noindent\textbf{Original atlas identifier:} P8; entry 48. \textbf{Field:} Public finance and taxation. \textbf{Original tractability label:} Medium.

\noindent\textbf{Classification:} Parameterized theoretical/empirical research agenda; not a certified unsolved theorem.

\subsection*{Definitions and assumptions}
Let jurisdictions $j=1,\ldots,J$ provide public goods $g_j\geq0$ at resource cost $C_j(g_j)$. Household $i$ chooses residence and consumes local services with spillovers represented by a specified matrix $A=(A_{jk})$; mobility incurs type-dependent costs. A constitutional policy $p$ assigns tax instruments, service authority, and grant rules to central and local governments. A model $m$ specifies voting/representation, tax competition, household sorting, and an equilibrium selection; decentralization need not implement a benevolent planner's allocation. Let $EV_i(p;m)$ be monetary equivalent variation, including taxes and mobility costs, and let $R_j(p;m)-C_j(p;m)$ be each government's net funds, counting intergovernmental transfers consistently so they cancel when aggregated. Fix welfare weights $w_i\geq0$ and fiscal values $\lambda_j$.

\subsection*{Formal research question}
\emph{Editorial formulation: the mathematical target below is not a theorem quoted from the references.}

For legal assignments $p\in\mathcal P$, identify and compare
\[
 W_m(p)=\mathbb E_m[w_iEV_i(p;m)]+\sum_{j=1}^{J}\lambda_j[R_j(p;m)-C_j(p;m)].
\]
Characterize which assignments and grant formulas are $\varepsilon$-optimal conditional on the estimated spillover matrix, mobility costs, information structure, and political mechanism. Determine whether boundary and grant reforms separately identify public-good spillovers, sorting, and strategic tax responses; give observational-equivalence examples or bounds when they do not. If public services have multidimensional quality, replace scalar $g_j$ by a defined service vector before comparing functions. Assess implementability in the selected political equilibrium rather than assuming the planner can directly set every local decision.

\subsection*{Interpretation and scope}
The social boundary must be fixed: a local government's gain from attracting taxpayers can be another jurisdiction's loss, while migration can also improve matching to services. Aggregate grants are transfers, but unequal marginal values of public funds can make their distribution relevant. The matrix $A$ need not be symmetric, and a central government may lack local preference information; both properties influence the optimal assignment. Boundary reforms often change electoral representation and tax bases alongside service provision, so exclusion restrictions require institutional support. The empirical extension estimates these mechanisms jointly and evaluates constitutional assignments under them. Tiebout mobility and classical decentralization arguments are conditional benchmarks, not claims that competition or decentralization always maximizes welfare.

\subsection*{Source audit and status}
All three original sources are identifiable. [S2]'s 1972 original is distinguished from the publisher's 2011 reprint. [S3] is explicitly political economy, so replacing it by a benevolent-government theorem would misstate the source. No universal decentralization open theorem is supplied by the atlas.

\subsection*{References}
\begin{enumerate}[label={[S\arabic*]},leftmargin=*,itemsep=0.6em]
\item Charles M. Tiebout (1956). A Pure Theory of Local Expenditures. Journal of Political Economy 64(5), 416--424. DOI: 10.1086/257839.
\url{https://www.journals.uchicago.edu/doi/10.1086/257839}
\newline\emph{Locator and support limits:} Article: mobility and local-public-good benchmark under restrictive assumptions; not general decentralization optimality.
\item Wallace E. Oates (1972). Fiscal Federalism. New York: Harcourt Brace Jovanovich. Reprinted by Edward Elgar, 2011, ISBN 9780857939944.
\url{https://www.e-elgar.com/shop/gbp/fiscal-federalism-9780857939944.html}
\newline\emph{Locator and support limits:} Publisher description and contents: fiscal assignment, grants, and decentralization. Publisher verifies 1972 original versus 2011 reprint.
\item Timothy Besley and Stephen Coate (2003). Centralized versus Decentralized Provision of Local Public Goods: A Political Economy Approach. Journal of Public Economics 87(12), 2611--2637. DOI: 10.1016/S0047-2727(02)00141-X.
\url{https://www.sciencedirect.com/science/article/pii/S004727270200141X}
\newline\emph{Locator and support limits:} Abstract and political-economy model: centralization trade-offs; endogenous mobility is a proposed extension.
\end{enumerate}

\subsection*{Common conventions: Source mathematical conventions}

Unless an entry states otherwise, agent and firm sets are finite. State and action spaces are measurable, strategies and outcome maps are measurable, probability laws are specified, and expectations exist for the targets under discussion. Infinite-horizon discount factors lie in $(0,1)$ and discounted objectives are finite. Norms, tolerances and complexity input sizes must be specified for computational comparisons. Constants or primitives that appear locally are inputs, not empirical facts asserted by this catalogue.

\textbf{Identification.} A statistical model $m\in\mathcal M$ implies an observable law $P_m$ and target $\theta(m)$. At law $P$, its identified set is
\[
\Theta(P;\mathcal M)=\{\theta(m):m\in\mathcal M,\ P_m=P\}.
\]
Point identification holds when this set is a singleton, or equivalently when $P_{m_1}=P_{m_2}$ implies $\theta(m_1)=\theta(m_2)$ for all compatible models. Sharp bounds mean bounds attained, or approached when the boundary is not attained, by compatible models. Writing this set is a definition, not a solution: the substantive task is to establish informative restrictions and derive or estimate the set. An unrestricted-confounding model may give only trivial bounds.

\textbf{Inference.} Identification, estimability and valid uncertainty quantification are separate questions. Fix $\alpha\in(0,1)$. A confidence procedure $C_n$ over a declared law class $\mathcal P$ has asymptotic uniform coverage at level $1-\alpha$ when
\[
\liminf_{n\to\infty}\inf_{P\in\mathcal P}\Pr_P\{\theta(P)\in C_n\}\geq1-\alpha.
\]
For set-identified targets the coverage event must be defined separately, for example coverage of the full identified set. If an entry uses identification-indexed models instead of uniquely indexed laws, the same coverage convention is applied over those models. Pointwise limits do not establish uniform validity, and asymptotic validity does not guarantee finite-sample precision. Sample sizes, dependence structures and weak-identification sequences are part of each application.

\textbf{Counterfactuals.} $Y_i(a)$ is a potential outcome under a specified intervention, including its timing, implementation and versions. It is not $Y_i$ conditional on voluntarily choosing $a$. A full intervention vector permits interference; shorthand scalar treatment notation does not silently impose no interference. Assignment, selection, attrition, anticipation and measurement must be specified. A claim about an instrument's local effect is not automatically a population policy effect.

\textbf{Economic equilibrium.} A counterfactual model includes best responses, constraints, feasibility, market clearing, state transitions and expectations consistent with the stated information. When several equilibria exist, either a declared selector is an input or the target is set-valued. Comparisons across policy regimes must use the stated selection convention. Reduced-form relationships are not presumed invariant after an equilibrium-changing intervention.

\textbf{Optimization and welfare.} Optimization is over the stated feasible set. If attainment is not established, $\sup$ or $\inf$ and $\varepsilon$-optimal choices replace an asserted maximizer or minimizer. For example, with value $V=\sup_{a\in\mathcal A}W(a)<\infty$, an $\varepsilon$-optimal set is $\{a\in\mathcal A:W(a)\geq V-\varepsilon\}$ for $\varepsilon>0$. Benefits, costs, risk and health or environmental outcomes can be aggregated only using declared units and conversion weights. Interpersonal utility weights, the treatment of fiscal transfers and foreign profits, discounting, and the behavioral welfare criterion are explicit inputs. A comparative-static derivative additionally requires existence and appropriate differentiability of the selected counterfactual.

\textbf{Sources.} Local labels [S1], [S2], and so forth restart in every question. Locators identify the relevant abstract, model, section or result at the level actually checked; a bibliographic or abstract-level verification is not represented as inspection of an entire proof. Working papers are distinguished from later publications and changing versions. Source summaries are paraphrased. All URL checks and editorial formulations refer to the audit date stated above.
% END ECONOMICS STATEMENT
\end{document}
